% ===== PRÉAMBULE CANONIQUE — à coller VERBATIM en tête de CHAQUE .tex =====
\documentclass[11pt,a4paper]{article}
\usepackage[utf8]{inputenc}\usepackage[T1]{fontenc}\usepackage{lmodern}
\usepackage[french]{babel}
\usepackage{amsmath,amssymb,mathtools}\usepackage{icomma}
\usepackage[version=4]{mhchem}          % \ce{...} : équations & formules
\usepackage{siunitx}\sisetup{locale=FR} % \qty{}{}, \si{}, \ang{}
\usepackage{chemfig}                    % \chemfig{...} : structures développées
\usepackage{xcolor}\usepackage{colortbl}
\usepackage{tikz}\usepackage{pgfplots}\pgfplotsset{compat=1.18}
\usepackage[most]{tcolorbox}\usepackage{enumitem}\usepackage{array,booktabs}
\usepackage[a4paper,margin=2.2cm]{geometry}
\usetikzlibrary{arrows.meta,calc,angles,quotes,positioning,patterns,backgrounds,shapes.geometric}
\definecolor{primary}{RGB}{222,49,99}\definecolor{primarydark}{RGB}{150,22,55}
\definecolor{defcol}{RGB}{0,114,178}\definecolor{methcol}{RGB}{52,92,148}
\definecolor{excol}{RGB}{0,135,90}\definecolor{attncol}{RGB}{200,40,55}
\tcbset{boxbase/.style={enhanced,breakable,boxrule=0.9pt,arc=2.5pt,
  left=3mm,right=3mm,top=2.3mm,bottom=2.3mm,fonttitle=\bfseries,coltitle=white}}
% --- boîtes TUTORIEL (noms EXACTS, ne pas renommer) ---
\newtcolorbox{cle}[1][]{boxbase,colback=defcol!6,colframe=defcol,
  title={\textsc{À retenir}\ifx\relax#1\relax\relax\else\textmd{ : \itshape #1}\fi}}
\newtcolorbox{methode}[1][]{boxbase,colback=methcol!7,colframe=methcol,sharp corners,
  title={\textsc{Méthode}\ifx\relax#1\relax\relax\else\textmd{ --- \itshape #1}\fi}}
\newtcolorbox{exemplebox}[1][]{boxbase,colback=excol!5,colframe=excol,
  title={\textsc{Exemple}\ifx\relax#1\relax\relax\else\textmd{ : \itshape #1}\fi}}
\newtcolorbox{piege}[1][]{boxbase,colback=attncol!6,colframe=attncol,
  title={\textsc{Piège}\ifx\relax#1\relax\relax\else\textmd{ : \itshape #1}\fi}}
\newtcolorbox{remarque}[1][]{boxbase,colback=black!4,colframe=black!45,
  title={\textsc{Remarque}\ifx\relax#1\relax\relax\else\textmd{ : \itshape #1}\fi}}
% --- boîtes EXERCICE / CORRIGÉ (noms EXACTS, auto-comptées) ---
\newtcolorbox[auto counter]{exo}[1][]{boxbase,colback=primary!4,colframe=primary,
  title={\textsc{Exercice}~\thetcbcounter\ifx\relax#1\relax\relax\else\textmd{ --- \itshape #1}\fi}}
\newtcolorbox[auto counter]{corrige}[1][]{boxbase,colback=excol!4,colframe=excol,
  title={\textsc{Corrigé}~\thetcbcounter\ifx\relax#1\relax\relax\else\textmd{ --- \itshape #1}\fi}}
\newcommand{\R}{\mathbb{R}}\newcommand{\N}{\mathbb{N}}\newcommand{\Z}{\mathbb{Z}}\newcommand{\ds}{\displaystyle}
% (un \usepackage{fancyhdr}+en-tête est possible ; il est ignoré côté web)
% =========================================================================
\begin{document}

\begin{center}{\large\bfseries\color{primarydark} Thème 4 — Niveau Difficile}\end{center}
\emph{Données (\si{\gram\per\mole}) :} \ce{NaCl}$=58{,}5$, \ce{CaCl2}$=111$.

\begin{exo}[trois sources d'ions chlorure]
Dans $\SI{1}{\litre}$ de solution, on dissout $\SI{0.05}{\mole}$ de \ce{NaCl},
$\SI{0.03}{\mole}$ de \ce{CaCl2} et $\SI{0.02}{\mole}$ de \ce{AlCl3}. Quelle est $[\ce{Cl-}]$~?
\end{exo}

\begin{exo}[mélange de deux nitrates]
On mélange $\SI{50}{\milli\litre}$ de \ce{Mg(NO3)2} à $\SI{0.010}{\mole\per\litre}$ et
$\SI{50}{\milli\litre}$ de \ce{Al(NO3)3} à $\SI{0.020}{\mole\per\litre}$. Le volume total est
$\SI{100}{\milli\litre}$.
\begin{enumerate}[label=\alph*)]
  \item Quelle est $[\ce{NO3-}]$ dans le mélange~?
  \item Quelles sont $[\ce{Mg^2+}]$ et $[\ce{Al^3+}]$~?
\end{enumerate}
\end{exo}

\begin{exo}[électroneutralité]
Une solution ne contient, comme ions, que \ce{Na+} à $\SI{0.10}{\mole\per\litre}$, \ce{Ca^2+} à
$\SI{0.05}{\mole\per\litre}$ (cations) et des ions chlorure \ce{Cl-} (seul anion). En écrivant que la
solution est électriquement neutre, détermine $[\ce{Cl-}]$.
\end{exo}

\begin{exo}[remonter à la masse de sel]
On veut préparer $\SI{250}{\milli\litre}$ d'une solution telle que $[\ce{Cl-}] = \SI{0.40}{\mole\per\litre}$,
en n'utilisant que du chlorure de calcium \ce{CaCl2}.
\begin{enumerate}[label=\alph*)]
  \item Quelle concentration $C$ en \ce{CaCl2} faut-il~?
  \item Quelle masse de \ce{CaCl2} faut-il dissoudre~?
\end{enumerate}
\end{exo}

\begin{exo}[toutes les concentrations d'ions]
On dissout $\SI{5.85}{\gram}$ de \ce{NaCl} et $\SI{11.1}{\gram}$ de \ce{CaCl2} dans de l'eau pour
obtenir $\SI{500}{\milli\litre}$ de solution.
\begin{enumerate}[label=\alph*)]
  \item Détermine $[\ce{Na+}]$ et $[\ce{Ca^2+}]$.
  \item Détermine $[\ce{Cl-}]$.
  \item Vérifie l'électroneutralité de la solution.
\end{enumerate}
\end{exo}

\end{document}
